% Supplementary Material of Paper A (separate file; not counted in the main-text page limit).
% Cross-references to the main text use xr: compile main.tex first, then this file, then main.tex again.
\documentclass[11pt]{article}
\usepackage[margin=2.5cm]{geometry}
\usepackage{times}
\usepackage[numbers,sort&compress]{natbib}
\usepackage{amsmath,amssymb,amsthm,booktabs,graphicx,array,longtable}
\newcommand{\widefig}[1]{\makebox[\linewidth][c]{\includegraphics[width=16cm]{#1}}}
\usepackage{xr-hyper}
\usepackage[colorlinks=true,allcolors=blue]{hyperref}
\externaldocument[][nocite]{build/main}[Paper_A_draft_EN.pdf]
\renewcommand{\thesection}{S\arabic{section}}
\renewcommand{\thefigure}{S\arabic{figure}}
\renewcommand{\thetable}{S\arabic{table}}
\newtheorem{proposition}{Proposition}
\renewcommand{\theproposition}{S\arabic{proposition}}
\renewcommand{\theequation}{S\arabic{equation}}
\title{Supplementary Material for ``When the site predicts the diagnosis: evaluation, bridging and within-site training for multi-site neuroimaging classifiers''}
\date{}
\pdfsuppressptexinfo=-1% do not copy the producer metadata of included figure files
\begin{document}
\maketitle
\noindent Section, figure and table numbers without the prefix S refer to the main text.

% Appendices A-D (written 2026-11-23 from docs/THEORY_A_PROOFS.md and the saved outputs; symbols as in Section 3).
\section{Proofs}\label{app:proofs}

Throughout, the model is Eq.~\eqref{eq:model}; the counts $a_s,t_s$ are fixed or independent of $u$ and $\varepsilon$, with $N_A,N_T>0$. The $\varepsilon_i$ are i.i.d.\ with mean $0$ and covariance $\Sigma_\varepsilon$ and independent of the $u_s$. $\Delta$ is a mean difference, not a causal effect. The AUC of a score $f$ is $P(f_{\mathrm{case}}>f_{\mathrm{ctrl}})+\tfrac12P(f_{\mathrm{case}}=f_{\mathrm{ctrl}})$, so a constant score has AUC $0.5$.

\subsection{Proposition~\ref{prop:pure}}

\emph{Part (a).} At a class-pure site $s$ every participant has $y_i=\gamma_s$. Under the second parameter value an observation has mean $\mu+(\Delta+\delta)\gamma_s+(u_s-\delta\gamma_s)=\mu+\Delta\gamma_s+u_s$, the same as under the first. Labels and the law of $\varepsilon$ are unchanged, so the training data have the same distribution under both values, and so does any estimator $T$. With $Z=T-\Delta$,
\[
\tfrac12\big(\mathbb{E}\|Z\|^2+\mathbb{E}\|Z-\delta\|^2\big)=\mathbb{E}\|Z-\delta/2\|^2+\|\delta\|^2/4\ \ge\ \|\delta\|^2/4,
\]
and the larger of the two risks is at least their average.

\emph{Part (b).} Let $\hat\Delta$ be the mean of the $k$ case-site means minus the mean of the $k$ control-site means, each site holding $m$ participants. Substituting Eq.~\eqref{eq:model},
\[
\hat\Delta=\Delta+\Big(\frac1k\sum_{\gamma_s=1}u_s-\frac1k\sum_{\gamma_s=0}u_s\Big)+(\bar\varepsilon_A-\bar\varepsilon_T).
\]
The common mean $m_u$ cancels in the bracket, so (A1$'$) suffices for unbiasedness. Under (A1) the bracket has covariance $(2/k)\Sigma_u$ and the noise term $(2/(km))\Sigma_\varepsilon$; the two are independent. The statement concerns the mean contrast, not the AUC of a particular classifier.

\subsection{Proposition~\ref{prop:L}}

Under (A0), $\bar x_A=\mu+\Delta+\sum_s(a_s/N_A)u_s+\bar\varepsilon_A$ and $\bar x_T=\mu+\sum_s(t_s/N_T)u_s+\bar\varepsilon_T$; subtracting gives the decomposition. Since $\sum_sc_s=1-1=0$, the mean $m_u$ cancels, and the contrast is unbiased under (A1$'$). Independent sites contribute $\sum_sc_s^2\Sigma_u=L\Sigma_u$, and cases and controls are disjoint, which gives the noise term. If site effects are correlated or heteroscedastic, the site term is $\sum_{s,r}c_sc_r\operatorname{Cov}(u_s,u_r)$; if (A1$'$) fails, a bias $\sum_sc_s\mathbb{E}u_s$ appears.

$L=0$ if and only if $c_s=0$ for every site, that is, $a_s/N_A=t_s/N_T$, which for a non-empty site means that its case fraction equals $N_A/(N_A+N_T)$.

For the symmetric designs, each class occupies $k$ originally pure sites with $m$ participants. With two-way bridging, $N_A=N_T=k(m+b)$, and $c_s=\pm(m-b)/(k(m+b))$ on case and control sites. With one-way bridging, $N_A=km$ and $N_T=k(m+b)$, and $c_s=\pm m/(k(m+b))$. The class-pure design is $b=0$. Squaring and summing over $2k$ sites gives Eq.~\eqref{eq:L-special}. In all three designs $c_s=+\alpha$ at the sites that contribute mainly cases and $-\alpha$ at the others, with $\alpha=\sqrt{L/(2k)}$, provided that $b\le m$ in the two-way design; for $b>m$ the signs swap.

\subsection{Corollary~\ref{cor:alloc}}

Add $b_j\ge0$ controls to case site $j$ and $d_j\ge0$ cases to control site $j$, with $U=\sum_jb_j$, $V=\sum_jd_j$ and $B=U+V$. The statements are for real-valued allocations.

\emph{Part (i).} For fixed $U$ and $V$ the totals $N_A=km+V$ and $N_T=km+U$ are fixed, and
\[
L=\sum_j\Big(\frac m{N_A}-\frac{b_j}{N_T}\Big)^2+\sum_j\Big(\frac{d_j}{N_A}-\frac m{N_T}\Big)^2
\]
is a separable, strictly convex function of the $b_j$ and of the $d_j$, each set under one linear constraint. The minimum is at equal allocation. With integer counts and $k$ not dividing $U$, the most even integer allocation is optimal by the same separable convexity.

\emph{Part (ii).} Write $u=U/k$, $v=V/k$, $t=B/k$ and $p=uv$. With equal allocation within each side, a case site has $c=\frac1k\big[\frac m{m+v}-\frac u{m+u}\big]=\frac1k\frac{m^2-uv}{(m+u)(m+v)}$ and a control site the negative, so
\[
L=\frac2k\,g(p)^2,\qquad g(p)=\frac{m^2-p}{m^2+mt+p}.
\]
Here $g'(p)=-(2m^2+mt)/(m^2+mt+p)^2<0$. If $B\le2km$, then $t\le2m$ and $p\le t^2/4\le m^2$, so $g\ge0$ and $g^2$ decreases in $p$; $p$ is largest at $u=v$. The condition cannot be dropped: for $k=1$, $m=2$, $B=5$, the split $(u,v)=(4,1)$ gives $uv=m^2$ and $L=0$, the even split gives $L=0.0247$, and the best near-even integer split $(3,2)$ gives $L=0.02$. All experimental cells have $b\le m$ and hence $B\le2km$.

\subsection{Proposition~\ref{prop:within}}

\emph{Part (a).} For a case $a$ and a control $t$ at the same site, $x_a-x_t=\Delta+\varepsilon_a-\varepsilon_t$ under (A0). A deterministic procedure whose input is only the set of these differences, with a fixed penalty and tuning rule, receives the same input after any shift $u_s\mapsto u_s+\eta_s$.

\emph{Part (b).} Assume $H>0$, that is, at least one training site holds both classes. For a scalar feature, regress $x$ on $y$ and one indicator per site. By the Frisch--Waugh--Lovell theorem the slope is $\sum_i(y_i-\bar y_{s(i)})x_i/\sum_i(y_i-\bar y_{s(i)})^2$. At site $s$, $\sum_{i\in s}(y_i-\bar y_s)^2=a_st_s/(a_s+t_s)=h_s$ and $\sum_{i\in s}(y_i-\bar y_s)x_i=h_sD_s$, which gives $\tilde\Delta=H^{-1}\sum_sh_sD_s$. A class-pure site contributes nothing to either sum, and its $D_s$ is not defined, so sums involving $D_s$ run over mixed sites ($a_st_s>0$). At each mixed site $D_s$ has covariance $(1/a_s+1/t_s)\Sigma_\varepsilon=\Sigma_\varepsilon/h_s$, and the $D_s$ are independent, so $\tilde\Delta$ is unbiased with covariance $\Sigma_\varepsilon/H$. If $H=0$ the slope is not defined.

L3 and the exact conditional logistic likelihood are different objectives. For a site with one case and two controls, with score differences $d_1$ and $d_2$, the conditional likelihood contributes $\log(1+e^{-d_1}+e^{-d_2})$ and L3 contributes $\log(1+e^{-d_1})+\log(1+e^{-d_2})=\log(1+e^{-d_1}+e^{-d_2}+e^{-d_1-d_2})$; the two coincide only for strata with one case and one control.

\subsection{Score gaps of linear scorers}

Section~\ref{sec:prop4} summarizes the following result, which was derived after the experiments.

\begin{proposition}[Leakage in score gaps]\label{prop:leak}
Fix a symmetric $Q\succeq0$ and consider the pooled linear scorer $w=Q(\bar x_A-\bar x_T)$ and the within-site scorer $\tilde w=Q\tilde\Delta$. Score a new case from training site $s$ and a new control from training site $r$, and let $G$ be the score gap (case minus control). Under (A0) and (A1), with new participants whose residuals are independent of the training sample given the shared site effects:
(a) for the pooled scorer, $\mathbb{E}\,G=\Delta^\top Q\Delta+(c_s-c_r)\operatorname{tr}(Q\Sigma_u)$;
(b) in the symmetric designs (with $b\le m$ under two-way bridging), a reverse-set pair has $c_s-c_r=-\sqrt{2L/k}$, so $\mathbb{E}\,G_{\mathrm{rev}}=\Delta^\top Q\Delta-\sqrt{2L/k}\,\operatorname{tr}(Q\Sigma_u)$;
(c) for a same-site pair ($s=r$) the pooled scorer has $\mathbb{E}\,G=\Delta^\top Q\Delta$;
(d) for the within-site scorer, $\mathbb{E}\,G=\Delta^\top Q\Delta$ for every pair.
\end{proposition}

\emph{Proof.} Assume (A0) and (A1). New participants share the site effect $u_s$ of their site with the training participants of that site; their residuals $\varepsilon$ are independent of the training sample. By Proposition~\ref{prop:L}, $w=Q\Delta+Q\sum_qc_q(u_q-m_u)+Q\bar e$ with $\bar e=\bar\varepsilon_A-\bar\varepsilon_T$, where $\sum_qc_q=0$ allows centring. A new case from site $s$ and a new control from site $r$ give $x_{\mathrm{case}}-x_{\mathrm{ctrl}}=\Delta+(u_s-m_u)-(u_r-m_u)+\varepsilon_{\mathrm{case}}-\varepsilon_{\mathrm{ctrl}}$. Taking the expectation of $G=w^\top(x_{\mathrm{case}}-x_{\mathrm{ctrl}})$, all cross terms vanish by independence except $\Delta^\top Q\Delta$ and
\[
\mathbb{E}\Big[\Big(\sum_qc_q(u_q-m_u)\Big)^{\!\top}Q\big((u_s-m_u)-(u_r-m_u)\big)\Big]=(c_s-c_r)\operatorname{tr}(Q\Sigma_u),
\]
which proves (a). Part (b) uses $c_s-c_r=-2\alpha=-\sqrt{2L/k}$ for a case from a control site and a control from a case site; in the two-way design this requires $b\le m$. Part (c) is $s=r$. For (d), $\tilde\Delta=\Delta+H^{-1}\sum_{s:a_st_s>0}h_s(\bar\varepsilon_{sA}-\bar\varepsilon_{sT})$ contains no site term, so only $\Delta^\top Q\Delta$ remains. These are statements about the mean score gap of linear scorers with a fixed $Q$; equal mean gaps do not imply equal AUCs, and the proposition says nothing about the AUC of regularized logistic regression or about how its penalty is selected.

\subsection{Within-site and cross-site targets}

Section~\ref{sec:eval} uses the following result.

\begin{proposition}[Two targets]\label{prop:targets}
For Gaussian $u$ and $\varepsilon$ and a fixed $w$ with $w^\top\Delta>0$ and $w^\top\Sigma_\varepsilon w>0$,
\begin{equation}
\mathrm{AUC}_W=\Phi\!\left(\frac{w^\top\Delta}{\sqrt{2w^\top\Sigma_\varepsilon w}}\right)\ \ge\
\mathrm{AUC}_B=\Phi\!\left(\frac{w^\top\Delta}{\sqrt{2w^\top(\Sigma_\varepsilon+\Sigma_u)w}}\right),\label{eq:targets}
\end{equation}
where $\mathrm{AUC}_W$ ranks a case and a control from the same new site and $\mathrm{AUC}_B$ a case and a control drawn from two independent new sites.
\end{proposition}

\emph{Proof.} Assume in addition that $u$ and $\varepsilon$ are Gaussian and that $w^\top\Sigma_\varepsilon w>0$. For a case and a control from the same new site, the score difference is Gaussian with mean $w^\top\Delta$ and variance $2w^\top\Sigma_\varepsilon w$, because the site effect cancels; for two independent new sites the variance is $2w^\top(\Sigma_\varepsilon+\Sigma_u)w$. Ties have probability zero, so each AUC equals $\Phi$ of the mean divided by the standard deviation. If $w^\top\Delta>0$, the larger variance gives the smaller AUC.

\subsection{Site-invariant representations under class-pure training}

Suppose every training site is class-pure, $Y=\gamma(S)$, and a representation $Z$ satisfies $P(Z\in A\mid S=s)=R(A)$ for all sites $s$ and measurable sets $A$. Then $P(Z\in A\mid Y=y)=\sum_{s:\gamma_s=y}R(A)P(S=s\mid Y=y)=R(A)$, so $Z$ is independent of $Y$, and any score computed from $Z$ alone has AUC $0.5$ on the training distribution. The argument needs exact marginal invariance, class-pure sites only, and scores that use $Z$ only; training with a finite adversarial weight guarantees none of these. It is a special case of the conflict between marginal alignment and label shift~\cite{zhao2019}.

\subsection{Monte Carlo check}

Table~\ref{tab:mc} compares the formulas with simulations on synthetic data ($d=3$, seed $0$).

\begin{table}[h]
\centering
\caption{Monte Carlo check of the formal results on synthetic data.}
\label{tab:mc}
\small
\resizebox{\linewidth}{!}{%
\begin{tabular}{lcc}
\toprule
Item & Simulated & Formula\\
\midrule
Prop.~\ref{prop:L}, $L$ for $k=2$, $m=20$, $b=5$ two-way & $0.36$ & $0.36$\\
Prop.~\ref{prop:L}, covariance, largest relative error ($20{,}000$ draws) & $0.009$ & $0$\\
Cor.~\ref{cor:alloc}, $k=3$, $m=10$, $B\in[0,60]$: $\min L-L_{\mathrm{equal}}$ & $7\times10^{-18}$ & $\ge0$\\
Prop.~\ref{prop:within}(b), fixed-effects slope against $\sum_sh_sD_s/H$ & $0.162732$ & $0.162732$\\
Prop.~\ref{prop:leak}(b), reverse-set gap, pooled scorer ($Q=I$, $40{,}000$ draws) & $-0.3668$ & $-0.3665$\\
Prop.~\ref{prop:leak}(c), same-site gap, pooled scorer & $0.2095$ & $0.21$\\
Prop.~\ref{prop:leak}(d), reverse-set gap, within-site scorer & $0.2140$ & $0.21$\\
Prop.~\ref{prop:targets}, $\mathrm{AUC}_W$ / $\mathrm{AUC}_B$ ($400{,}000$ draws) & $0.6303$ / $0.6198$ & $0.6292$ / $0.6201$\\
\bottomrule
\end{tabular}}
\end{table}

% Moved from Section 3 to save main-text pages (appendices are not counted by PR).
\subsection{A synthetic illustration of Proposition~\ref{prop:leak}}\label{sec:synthetic}

Table~\ref{tab:synthetic} evaluates Proposition~\ref{prop:leak}(b) for the plug-in scorer ($Q=I$) in a three-dimensional example with $\|\Delta\|^2=0.21$ and $\operatorname{tr}\Sigma_u=0.96$, using the design of our ABIDE-I experiments ($k=2$, $m=20$). Without bridges, the expected reverse-set gap is $-0.75$, more than three times the size of the true signal and of opposite sign. Two-way bridging with $b=5$ shrinks it to $-0.37$ and with $b=10$ to $-0.11$; one-way bridging with the same number of bridges per site removes less. A Monte Carlo check with $40{,}000$ replicates of the two-way $b=5$ design gives $-0.367$ for the pooled scorer (formula $-0.366$), $+0.214$ for the within-site scorer (formula $+0.21$), and $+0.210$ for a same-site pair scored by the pooled scorer.

\begin{table}[t]
\centering
\caption{Expected reverse-set score gap $\mathbb{E}\,G_{\mathrm{rev}}$ of the pooled plug-in scorer in the synthetic example ($k=2$, $m=20$, $\|\Delta\|^2=0.21$). The within-site scorer has $\mathbb{E}\,G=0.21$ in every design that contains bridges.}
\label{tab:synthetic}
\begin{tabular}{lcccc}
\toprule
& $b=0$ & $b=2$ & $b=5$ & $b=10$\\
\midrule
One-way: $L$ & 1.000 & 0.826 & 0.640 & 0.444\\
One-way: $\mathbb{E}\,G_{\mathrm{rev}}$ & $-0.751$ & $-0.663$ & $-0.559$ & $-0.431$\\
Two-way: $L$ & 1.000 & 0.669 & 0.360 & 0.111\\
Two-way: $\mathbb{E}\,G_{\mathrm{rev}}$ & $-0.751$ & $-0.576$ & $-0.366$ & $-0.110$\\
\bottomrule
\end{tabular}
\end{table}


\section{Pre-specified criteria}\label{app:criteria}

Each experiment was specified in writing before it was run, including the cells, the repeats, the primary metric and the criteria that would count as success. Table~\ref{tab:criteria} lists every criterion in the order in which the criteria were fixed, with its outcome, including criteria for methods that are not reported in the main text. Three analyses in the main text were not pre-specified and are labelled as exploratory or post hoc where they appear: the rank correlation between $L$ and the reverse-set AUC over cells that had been inspected before Proposition~\ref{prop:L} was written, the AUC version of the difference-in-differences for KUL\_3, and the comparison of the two settings for the published methods (whose own criteria N1 and N2 were fixed before it was run). Proposition~\ref{prop:leak} was derived after the experiments. In the deep pilot, the learning rate of the autoencoder was lowered from $10^{-3}$ to $10^{-4}$ after the training loss diverged; this change was judged on training losses only.

{\footnotesize\setlength{\tabcolsep}{3pt}
\begin{longtable}{>{\raggedright}p{0.19\linewidth}>{\raggedright}p{0.38\linewidth}>{\raggedright}p{0.2\linewidth}>{\raggedright\arraybackslash}p{0.13\linewidth}}
\caption{Pre-specified criteria and outcomes, in the order in which they were fixed. ABIDE-I with CC200 unless stated otherwise; values are means over repeats.}\label{tab:criteria}\\
\toprule
Analysis & Criterion & Result & Met\\
\midrule
\endfirsthead
\toprule
Analysis & Criterion & Result & Met\\
\midrule
\endhead
\bottomrule
\endfoot
Cost of purity & P1: M $-$ P $\ge0.05$ at $k=1$ & $0.057$ & yes\tabularnewline
 & P2: P higher and cost lower at $k=6$ than at $k=1$ & $+0.017$; $0.057\to0.015$ & yes\tabularnewline
 & P3: from $n=20$ to $80$ at $k=2$, P gains less than M & $0.034$ vs $0.080$ & yes\tabularnewline
Evaluation & V1: $E_{\mathrm{subj}}-$ truth $\ge0.10$ at $k=2$ & $+0.365$ & yes\tabularnewline
 & V2: LSPO closer to truth than $E_{\mathrm{subj}}$ at $k=2,4,6$ & MAE $0.16/0.11/0.09$ vs $0.37/0.27/0.19$ & yes\tabularnewline
 & M1: LSPO tuning does not lower test AUC and raises reverse-set AUC ($k=2$; $k=4$) & $+0.001/+0.004$; $-0.000/+0.003$ & $k=2$ only\tabularnewline
 & M2: NAP raises test and reverse-set AUC at $k=2$ and $4$ & $-0.001$, $-0.000$ & no\tabularnewline
Explicit site and diagnosis correction before LR (C2; not in main text) & R1: two-way $b=5$ recovers $\ge50\%$ of the cost and reverse-set AUC $\ge0.40$ & $84\%$; $0.07$ & no\tabularnewline
 & R2: C2 $-$ C0 test AUC $>0$ at $b=2$ and $5$ & mixed signs & no\tabularnewline
 & R3: C2 (one-way, $b=5$, $+10$ subjects) $\ge$ P with $n=80$ ($+80$) & $+0.002$ & yes\tabularnewline
Same-class control & D1: two-way $b=5$, C0 $-$ S0 $>0$ in test and reverse-set AUC & $+0.044$, $+0.24$ ($k=2$) & yes\tabularnewline
 & D2: same, one-way & $+0.015$, $+0.10$ & yes\tabularnewline
Robustness (QC subset, HO, BrainNetCNN) & Q1: M $-$ P $\ge0.03$ at $k=1$, lower cost and higher reverse-set AUC at $k=6$; Q2 as V1/V2; Q3 as D1 & Table~\ref{tab:robust} & yes, all three\tabularnewline
Within-site model on mixed data (exploratory precursor) & N1: L3 $-$ L0 $\ge+0.01$ with $\ge70\%$ positive repeats & $+0.008$, $78\%$ & no\tabularnewline
 & N2: $|$L3 $-$ L1$|<0.005$ & $0.006$ & no\tabularnewline
Within-site estimator (CC200 and HO) & E1: in all four $b=5$ cells, L3 reverse-set AUC $\ge0.50$ and above C0 & $0.59$--$0.65$ & yes\tabularnewline
 & E2: two-way $b=5$, L3 $-$ C0 test AUC $\ge-0.01$ & $+0.005$ to $+0.009$ & yes\tabularnewline
Design rules (CC200 and HO) & G1: even minus concentrated allocation $>0$, two-way $k=4$ & $+0.040$, $+0.037$ & yes\tabularnewline
 & G2: Spearman$(L$, C0 reverse-set AUC$)\le-0.8$ & $-0.997$, $-0.991$ & yes\tabularnewline
 & G3: ComBat null test, reverse-set AUC in $[0.40,0.60]$ & EB: $0.000$ (4 of 4 groups fail); without EB: 3 of 4 pass & no\tabularnewline
 & G4: $|$L1 $-$ L3$|<0.005$ in test AUC & $0.003$--$0.004$ & yes\tabularnewline
ABIDE-II replication & R1--R5 as P1/Q1, V1/V2, D1, E1/E2, G2 & Table~\ref{tab:replication} & R3 not assessable; others yes\tabularnewline
Group DRO & G0: gradient ERM within $0.01$ of sklearn ERM & $0.001$--$0.003$ & yes\tabularnewline
 & H1: Group DRO $-$ ERM reverse-set AUC $<0.05$ at $k=2$ and $4$ & $-0.001$, $+0.000$ & yes\tabularnewline
KUL\_3 (ABIDE-II) & K1: DiD in sensitivity $\ge0.10$ & $+0.59$ & yes\tabularnewline
Mixed objective, tuned by within-site AUC & W1: reverse-set AUC $\ge$ L3 $-0.02$ in all four $b=5$ cells & $-0.05$ to $-0.18$ & no\tabularnewline
 & W2: test AUC $\ge\max($C0, L3$)-0.005$ in all four cells & 3 of 4 cells & no\tabularnewline
 & W3: tuning by subject folds lowers reverse-set AUC by $\ge0.05$ & $0.146$ & yes\tabularnewline
Deep pilot ($k=2$, $b=5$ two-way, repeats $0$--$9$; selects the candidate for Section~\ref{app:deep}) & T1: reverse-set AUC $\ge0.55$ and test AUC $\ge$ deep ERM $-0.01$, for each of three candidates & encoder on training data, L3 head: $0.49$; encoder on training and external data, L3 head: $0.59$; end-to-end pair loss: $0.49$ & second candidate only\tabularnewline
 & T2: for a candidate meeting T1, L3 head $-$ logistic head reverse-set AUC $\ge0.10$ & $+0.28$ & yes\tabularnewline
Deep within-site head (Section~\ref{app:deep}) & D1: E-L3 $-$ E-CE reverse-set AUC $\ge0.10$ in each cell; test difference, averaged over the four cells, $\ge-0.01$ & $0.07$--$0.24$ (one cell $<0.10$); $+0.000$ & no\tabularnewline
 & D2: E-L3 reverse-set AUC $\ge0.55$ and $\ge$ best deep baseline; test $\ge$ D-ERM $-0.01$ & $0.59$ vs $0.39$; $+0.004$ & yes\tabularnewline
ADHD-200 & H1: WashU DiD $\ge0.10$ & $+0.30$ & yes\tabularnewline
 & R1: M $-$ P $\ge0.03$ & $0.008$ & no\tabularnewline
 & R2: $E_{\mathrm{subj}}-$ truth $\ge0.10$ & $+0.46$ & yes\tabularnewline
 & R3: C0 $-$ S0 $>0$ in test and reverse-set AUC & $+0.013$, $+0.076$ & yes\tabularnewline
 & R4: L3 reverse-set AUC $\ge0.50$, above C0; test $\ge-0.01$ & $0.56$, $0.52$ & yes\tabularnewline
 & R1 trend, LSPO, R5 & only $k=1$ & not assessable\tabularnewline
One-class models & O-a: O1 reverse-set AUC $<0.5$ in B0 & $0.52$ & no\tabularnewline
 & O-b: O3 reverse-set AUC $\ge0.5$ and above O1, one-way $b=5$ & $0.25$ vs $0.50$ & no\tabularnewline
 & O3-eq: reverse-set AUC $\ge0.45$ in all four cells & $0.50$--$0.54$ & yes\tabularnewline
Same institution & S-a: P\_same reverse-set AUC above P\_cross & $0.55$ vs $0.20$ & yes\tabularnewline
Published methods (per method) & G1: B0 reverse-set AUC $<0.5$ at $k=2$ and $4$ & $k=4$: $0.52$--$0.58$ & no (all)\tabularnewline
 & G2: $E_{\mathrm{subj}}-$ truth $\ge0.10$ & Parisot $+0.41$, $+0.40$; MIDA $-0.27$ & Parisot only\tabularnewline
 & G3: C0 $-$ B0 reverse-set AUC $>0$ at $k=2$ and $4$ & $+0.05$ to $+0.58$ & yes (all)\tabularnewline
 & N1 (post hoc): MIDA without reverse set, B0 reverse-set AUC $<0.5$ & $0.35$, $0.40$ & yes\tabularnewline
 & N2 (post hoc): graph network without reverse set, B0 $k=4$ & $0.22$ & yes\tabularnewline
Partial confounding (Section~\ref{app:partial}) & PC1: pooled reverse-set AUC falls with $\delta$; drop at $\delta=0.3$ $\ge0.05$ in $\ge80\%$ of repeats & $0.287$, $100\%$ & yes\tabularnewline
 & PC2: drop at $\delta=0.2$ $\ge0.03$ & $0.184$ & yes\tabularnewline
 & PC3: at $\delta=0.4$, L3 drop $\le$ half the pooled drop; L3 $-$ pooled test AUC $\ge-0.01$ for $\delta\le0.3$ & $0.054$ vs $0.391$; $\ge+0.001$ & yes\tabularnewline
 & PC4, PC5: weights, L1, centring and ComBat against L3; training-data estimates & Table~\ref{tab:partial}, Fig.~\ref{fig:partial}c & descriptive\tabularnewline
 & PC6: at $\delta=0.4$, class-conditional projection keeps $\ge$ half the pooled drop, and pooled-class projection loses $\ge0.02$ test AUC relative to the pooled model beyond their difference at $\delta=0$, each in $\ge80\%$ of repeats & $0.376$ vs $0.391$, $100\%$; $-0.023$, $70\%$ & first only\tabularnewline
\end{longtable}
}


% Supplementary Section S3 -- additional results and method details moved from the main text (2026-11-23).
\section{Additional results}\label{app:more}

Unless stated otherwise, results are for ABIDE-I with the CC200 atlas, as means over $50$ repeats.

\subsection{Main designs and methods}

\begin{table}[h]
\centering
\caption{Training data and training objective of the main designs and methods, for $k=2$, $m=20$ and $b=5$ (cases/controls contributed by each training site). Sites 1 and 2 are the case sites and sites 3 and 4 the control sites of the class-pure design. Pooled LR: $\ell_2$-penalized logistic regression on all training participants.}
\label{tab:designs}
\setlength{\tabcolsep}{3pt}
\makebox[\linewidth][c]{\begin{minipage}{15cm}\centering
\begin{tabular}{>{\raggedright}p{0.3\linewidth}cccccl}
\toprule
Name & Site 1 & Site 2 & Site 3 & Site 4 & Total & Objective\\
\midrule
B0, class-pure pooled LR (design P) & 20/0 & 20/0 & 0/20 & 0/20 & 80 & pooled LR\\
M, mixed-site pooled LR & 10/10 & 10/10 & 10/10 & 10/10 & 80 & pooled LR\\
S0, same-class additions & 25/0 & 25/0 & 0/25 & 0/25 & 100 & pooled LR\\
C0, bridged pooled LR & 20/5 & 20/5 & 5/20 & 5/20 & 100 & pooled LR\\
L3, bridged within-site LR & 20/5 & 20/5 & 5/20 & 5/20 & 100 & within-site pairs\\
\bottomrule
\end{tabular}
\end{minipage}}
\end{table}

\subsection{Results by number of sites}

Table~\ref{tab:kcurve} gives the numbers behind Figs.~\ref{fig:problem} and~\ref{fig:evaluation}a--c.

\begin{table}[h]
\centering
\caption{ABIDE-I, CC200, $n=40$ per class. Within-site AUC on held-out test sites for the pure (P) and mixed (M) designs, reverse-set AUC of P, and two estimates of the within-site AUC of P computed from the training data: subject-level five-fold cross-validation ($E_{\mathrm{subj}}$) and leave-site-pair-out ($E_{\mathrm{LSPO}}$). MAE: mean absolute error against the test-site value across repeats.}
\label{tab:kcurve}
\begin{tabular}{lcccc}
\toprule
& $k=1$ & $k=2$ & $k=4$ & $k=6$\\
\midrule
P, within-site AUC & 0.578 & 0.575 & 0.581 & 0.596\\
M, within-site AUC & 0.636 & 0.627 & 0.609 & 0.611\\
Cost M$-$P & 0.057 & 0.053 & 0.028 & 0.015\\
P, reverse-set AUC & 0.029 & 0.120 & 0.269 & 0.365\\
\midrule
$E_{\mathrm{subj}}$ (MAE) & n.a. & 0.940 (0.365) & 0.851 (0.269) & 0.785 (0.189)\\
$E_{\mathrm{LSPO}}$ (MAE) & n.a. & 0.523 (0.161) & 0.584 (0.115) & 0.592 (0.090)\\
\bottomrule
\end{tabular}
\end{table}

\subsection{Site identification and site information in diagnostic scores}

\emph{Methods.} A multinomial logistic regression predicted the site from the features, with nested five-fold cross-validation stratified by site and five seeds. Its balanced accuracy was compared with chance, $1/S$, and with $100$ runs on permuted site labels that reuse the penalty selected most often on the real labels, an approximate rather than exact permutation reference. It was repeated with controls only and with cases only. For trained B0, C0 and L3 models, the same-diagnosis site AUC is the AUC with which the diagnostic score separates unused participants of one diagnosis at case sites from those at control sites.

\emph{Results.} The site was identified with a balanced accuracy of $0.763$ among the $17$ ABIDE-I sites (chance $0.059$), $0.722$ among the $13$ ABIDE-II sites (chance $0.077$) and $0.919$ among the six ADHD-200 sites (chance $0.167$), above every permuted run (Fig.~\ref{fig:sitecode}a). Accuracy stayed far above chance with controls only ($0.683$, $0.548$, $0.911$) or cases only ($0.617$, $0.490$, $0.903$). The same-diagnosis site AUC of the class-pure model was $0.92$ at $k=2$ and $0.81$ at $k=4$; two-way bridging lowered it to $0.81$ and $0.64$, and the within-site estimator to $0.49$--$0.52$ (Fig.~\ref{fig:sitecode}b). An AUC near $0.5$ means that the score no longer ranks participants of the same diagnosis by site; it does not show that the score is independent of the site.

\begin{figure}[tbp]
\centering
\widefig{figures/fig_sitecode.pdf}
\caption{Sites are easy to identify from connectivity, and the score of a pooled model trained on class-pure sites encodes them. (a)~Balanced accuracy of a multinomial logistic regression that predicts the site from the connectivity features, using all participants (grey), controls only (aqua) or cases only (gold); bars are means over five seeds of nested five-fold cross-validation, black dots are $100$ runs with permuted site labels, and dashed lines mark chance, $1/S$. (b)~Same-diagnosis site AUC: within one diagnosis, the AUC with which the diagnostic score separates participants from case sites and from control sites ($0.5$: the score does not rank them by site). Shades of red: pooled models; shades of blue: the within-site estimator L3 ($b=5$). Dots are single repeats ($50$ per condition) averaged over the two diagnoses, large markers are means (ABIDE-I, CC200).}
\label{fig:sitecode}
\end{figure}

\subsection{Scores of individual participants}

The class-pure model ranks controls held out at case sites above the cases at new sites (median rank $0.74$ against $0.52$), and cases held out at control sites below the controls at new sites ($0.24$ against $0.43$; Fig.~\ref{fig:anatomy}). Two-way bridging moves both reverse-set groups towards the middle under the pooled model ($0.64$ and $0.39$), and the within-site estimator puts them in the same order as participants at new sites ($0.46$ for controls and $0.65$ for cases).

\begin{figure}[tbp]
\centering
\widefig{figures/fig_anatomy.pdf}
\caption{The site shortcut ranks subjects by their site ($k=2$, ABIDE-I, CC200). Within each repeat, the scores of all evaluated subjects were converted to ranks between $0$ (most control-like) and $1$ (most case-like), and the ranks were pooled over $50$ repeats. Filled densities: cases (gold) and controls (aqua) at held-out test sites. Dashed densities: the reverse set, that is, controls held out at originally case-pure sites and cases held out at originally control-pure sites. Black ticks mark medians. (a)~The class-pure model scores reverse-set controls as more case-like than the cases at new sites. (b)~Two-way bridging with $b=5$ weakens this inverted ordering under the pooled model but does not remove it. (c)~The within-site estimator orders the reverse set in the same way as subjects at new sites.}
\label{fig:anatomy}
\end{figure}

\subsection{Methods that use only class-pure data}

\emph{Methods.} Tuning by LSPO instead of subject-level folds; removal of the subspace spanned by site means (nuisance attribute projection, NAP); per-site standardization; and group distributionally robust optimization (Group DRO) with one group per site~\cite{sagawa2020}, compared with ERM trained by the same gradient optimizer.

\emph{Results.} None of these methods materially reduced the measured shortcut (Fig.~\ref{fig:pureonly}, Table~\ref{tab:methods}). LSPO tuning and NAP changed the within-site AUC by less than $0.002$ and the reverse-set AUC by less than $0.005$ at every $k$. Per-site standardization produced an AUC of exactly $0.500$ on every metric in every repeat: when each site holds one class, subtracting the site mean also subtracts the class mean. Group DRO differed from ERM by $-0.001$ ($k=2$) and $+0.000$ ($k=4$) in reverse-set AUC, also with a larger step size. The site shortcut already drove every group's training loss to $0.01$--$0.04$, so the group weights stayed close to uniform (largest weight $0.28$ against $0.25$ at $k=2$, and $0.17$ against $0.125$ at $k=4$).

\begin{figure}[tbp]
\centering
\widefig{figures/fig_pureonly.pdf}
\caption{Methods that use only class-pure data leave the shortcut in place (ABIDE-I, CC200, $n=40$). (a, b)~Paired differences from the pure baseline B0 in within-site AUC and in reverse-set AUC, for $k=2$ (circles) and $k=4$ (squares). Group DRO is compared with ERM trained by the same gradient optimizer. Methods on class-pure data are drawn in red; two-way bridging with $b=5$ (blue) is shown for scale. Per-site standardization is omitted because all of its AUCs equal $0.500$. (c)~Final Group DRO weight of each group (site) in each repeat; blue lines mark uniform weights $1/(2k)$. (d)~Final training loss of each group under Group DRO, on a logarithmic scale; black lines are medians and the dashed line is the loss $\log 2$ of a constant score. Dots are single repeats ($50$ per condition), large markers are means.}
\label{fig:pureonly}
\end{figure}

\begin{table}[h]
\centering
\caption{ABIDE-I, CC200, $k=2$, $n=40$. Methods that use only class-pure data, compared with the pure baseline B0. Paired differences from B0 are means over $50$ repeats. $^{\dagger}$Compared with ERM trained by the same gradient optimizer.}
\label{tab:methods}
\begin{tabular}{lcc}
\toprule
Method & $\Delta$ within-site AUC & $\Delta$ reverse-set AUC\\
\midrule
LSPO tuning & $+0.001$ & $+0.004$\\
NAP (site-mean subspace removed) & $-0.001$ & $+0.002$\\
Per-site standardization & all AUCs $=0.500$ & --\\
Group DRO (group $=$ site)$^{\dagger}$ & $+0.000$ & $-0.001$\\
\bottomrule
\end{tabular}
\end{table}

\subsection{One-class models}

\emph{Methods.} O1 is a Ledoit--Wolf Mahalanobis distance to the mean of the training controls in a $30$-component principal subspace; O2 is a one-class support vector machine; O3 is the O1 distance measured from the mean of the controls at the participant's own site, with half of the controls at each test site used as its reference and excluded from evaluation. A variant, O3-eq, uses five reference controls at every training site.

\emph{Results.} We had predicted, before running them, that O1 would rank the reverse set below chance and that O3 would not. Neither prediction held. O1 and O2 reached within-site AUCs of only $0.50$ and $0.55$ and reverse-set AUCs near $0.5$: they learned neither the site nor the diagnosis. O3 had a within-site AUC of $0.51$, and its reverse-set AUC fell to $0.18$--$0.31$ in the bridged designs, because a case site had only five bridging controls to estimate its reference whereas a control site had twenty. When every training site used five reference controls (O3-eq), the reverse-set AUC rose to $0.50$--$0.54$ in all four cells, as a check specified before it was run required. One-class training thus avoided the shortcut only by learning no diagnostic signal, and normative calibration requires enough controls at every site.

\subsection{Same-institution designs}

\emph{Methods.} LEUVEN, UCLA and UM each contributed two collections. In each repeat, one collection of each institution supplied only cases and the other only controls ($13$ per collection). This design was compared with $13$ per site from three other case-only and three other control-only sites, and with both classes drawn from all six collections. Test sites were drawn from the remaining sites.

\emph{Results.} With collections from the same institution, the reverse-set AUC was $0.55$ instead of $0.20$, higher in all $50$ repeats, and the within-site AUC on new sites was $0.64$ instead of $0.57$, $0.02$ below a design in which both collections held both classes (Fig.~\ref{fig:sameinst}). The benefit depended on how similar the two collections were: LEUVEN's two collections could be told apart with an AUC of $0.98$ and its reverse-set AUC was $0.27$; for UM these values were $0.74$ and $0.59$, and for UCLA $0.53$ and $0.68$. With three institutions this pattern is descriptive. It is consistent with the theory once site effects are allowed to be correlated: the site term becomes $\sum_{s,t}c_sc_t\operatorname{Cov}(u_s,u_t)$, and two collections that share a scanner and enter with opposite classes contribute a negative cross term.

\begin{figure}[tbp]
\centering
\widefig{figures/fig_sameinst.pdf}
\caption{Class-pure collections taken from the same institution carry less site shortcut than collections from different institutions (ABIDE-I, CC200; $39$ participants per class from three institutions, LEUVEN, UCLA and UM, or from six other sites). (a)~Within-site AUC on new sites and (b)~reverse-set AUC when cases and controls come from different institutions (red), from two collections of the same institution (blue), or from both collections with both classes (grey; no reverse set). Dots are single repeats ($50$), large markers are means. (c)~For each institution, the AUC with which its two collections can be told apart from their connectivity, against the reverse-set AUC of its participants in the same-institution design; three points, shown descriptively.}
\label{fig:sameinst}
\end{figure}

\subsection{Mixing the within-site and pooled objectives}

We mixed the two objectives as $(1-\lambda)\,\ell_{\mathrm{pair}}+\lambda\,\ell_{\mathrm{CE}}$ and traced $\lambda$ from $0$ (L3) to $1$ (C0) on the same training data. At $k=2$, $b=5$, one-way, moving from $\lambda=0$ to $\lambda=0.1$ lowered the reverse-set AUC from $0.603$ to $0.328$ while raising the within-site AUC by only $0.016$, and larger $\lambda$ changed little (Fig.~\ref{fig:lambda}a,b). When $\lambda$ was selected by the within-site pair AUC in inner folds, $\lambda=0$ was the most frequent choice but was chosen in only $41\%$ of repeats, and the selected models had reverse-set AUCs $0.05$--$0.18$ below L3 in the four $b=5$ cells. When $\lambda$ was selected by subject-level folds, $\lambda=0$ was chosen in $1\%$ of repeats (Fig.~\ref{fig:lambda}c). The pre-specified criteria W1 and W2 were not met (Table~\ref{tab:criteria}).

\begin{figure}[tbp]
\centering
\widefig{figures/fig_lambda.pdf}
\caption{Admitting even a small weight of the pooled cross-entropy brings the shortcut back, and tuning does not reliably avoid it (ABIDE-I, CC200, $b=5$). The objective is $(1-\lambda)\,\ell_{\mathrm{pair}}+\lambda\,\ell_{\mathrm{CE}}$, so $\lambda=0$ is L3 and $\lambda=1$ is C0; marker colour runs from blue ($\lambda=0$) to red ($\lambda=1$). (a)~Reverse-set AUC and (b)~within-site AUC against $\lambda$, means over $50$ repeats; solid lines $k=2$, dashed lines $k=4$, light grey one-way and dark grey two-way bridging. (c)~Distribution of the $\lambda$ selected in inner folds, pooled over the four cells, when selection maximizes the within-site pair AUC (blue) or the subject-level AUC (red).}
\label{fig:lambda}
\end{figure}

\subsection{Natural class-pure collections}

For KUL\_3, the AUC of its cases against mixed-site controls fell from $0.754$ under subject-level cross-validation to $0.430$ under leave-one-site-out evaluation, while that of mixed-site cases fell from $0.701$ to $0.677$, so that $\mathrm{DiD}=+0.30$ (mean over repeats $+0.299$ [$0.240$, $0.346$]); this AUC analysis is exploratory, added after the pre-specified analysis of sensitivity. For WashU the AUCs fell from $0.695$ to $0.305$ and from $0.591$ to $0.505$ ($+0.304$ [$0.276$, $0.343$], pre-specified), and $+0.29$ to $+0.33$ for males only, ages $7$--$13$, or separate Peking collections. For Pittsburgh ($92$ controls and $4$ cases), they fell from $0.802$ to $0.205$ and from $0.588$ to $0.495$ ($+0.50$ [$0.48$, $0.52$]).

\subsection{Published multi-site methods: results and implementation}

Figure~\ref{fig:published} shows the results summarized in Section~\ref{sec:res-pureonly}. In the original transductive setting, the reverse set entered training without labels and both methods looked better (reverse-set AUC up to $0.58$), plausibly because these participants acted as unlabeled bridges; this post hoc comparison also changes the pool size, the feasible MIDA dimension and the scoring.

\begin{figure}[!htbp]
\centering
\widefig{figures/fig_published.pdf}
\caption{Two published multi-site methods on class-pure training data (ABIDE-I, $n=40$; population graph network with the Harvard-Oxford atlas, MIDA pipeline and logistic regression with CC200). Red: class-pure design B0 with the reverse set kept out of training, as in a real class-pure collection. Grey, hollow: B0 in the original transductive setting, in which the reverse set entered training without labels. Blue: two-way bridging C0 with $b=5$. (a, b)~Reverse-set AUC at $k=2$ and $k=4$; the shaded region lies below chance. (c)~Within-site AUC on new sites at $k=2$; black lines mark the same method trained on mixed sites (M). (d)~Subject-level five-fold cross-validation of B0 at $k=2$ (red) against the within-site AUC on new sites (blue); bold numbers are the mean differences. Logistic regression is inductive. Numbers above the dots are means; dots are single repeats ($50$).}
\label{fig:published}
\end{figure}

We ported the TensorFlow code of the population graph network~\cite{parisot2018} line by line to PyTorch, fitted its recursive feature elimination on the training participants only, and used the Harvard-Oxford atlas of the original paper. We ran the authors' code of the MIDA pipeline~\cite{kunda2023}. Before the main runs, each method reproduced its published ABIDE-I accuracy within three percentage points, a margin fixed in advance ($69.2\%$ against $70.4\%$, and $71.9\%$ against $73.0\%$). The MIDA pipeline departed from the original in two respects. It received only age and sex as phenotypes, because full-scale IQ carries diagnostic information and eye status during the scan is set by the acquisition protocol and therefore identifies the site. Its MIDA dimension was searched only over values not larger than the number of participants in the fit; $h=300$ was unavailable in $32$ of $50$ B0 repeats when the reverse set was kept out of training. To score the reverse set after training, we projected it with the learned MIDA kernel, centred with the column means of the fitted kernel, and attached its nodes to the trained graph by the original edge rules followed by one forward pass; this out-of-sample scoring is our extension of the authors' code. Full results are in Table~\ref{tab:published}.

\subsection{Replication criteria}

Table~\ref{tab:replication} lists the pre-specified replication criteria of Section~\ref{sec:res-replication} (results in Fig.~\ref{fig:replication}).

\begin{table}[h]
\centering
\caption{Pre-specified replication criteria on ABIDE-II (Gordon atlas, $n=20$) and ADHD-200 (CC200, $n=20$, $k=1$ only). n.a.: not assessable in this cohort.}
\label{tab:replication}
\makebox[\linewidth][c]{\begin{minipage}{15cm}\centering
\begin{tabular}{p{0.34\linewidth}p{0.27\linewidth}p{0.27\linewidth}}
\toprule
Criterion & ABIDE-II & ADHD-200\\
\midrule
R1: cost M$-$P $\ge0.03$ at the smallest $k$, and smaller cost with more sites & met ($0.050\to0.000$) & first half not met ($0.008$); trend n.a.\\
R2: $E_{\mathrm{subj}}$ overestimates by $\ge0.10$ & met ($+0.38$) & met ($+0.46$); LSPO n.a.\\
R3: C0 $>$ S0 in test and reverse-set AUC & n.a. ($<20$ repeats) & met ($+0.013$, $+0.076$)\\
R4: L3 reverse-set AUC $\ge0.5$ and above C0; test cost $\ge-0.01$ (two-way) & met ($0.63$ vs.\ $0.39$; $-0.006$) & met ($0.52$ vs.\ $0.08$; $+0.004$)\\
R5: rank correlation of $L$ with C0 reverse-set AUC & met ($-0.970$, 11 cells) & n.a. (5 cells)\\
\bottomrule
\end{tabular}
\end{minipage}}
\end{table}

\subsection{Within-site estimation across bridge numbers}

Fig.~\ref{fig:tradeoff} shows the within-site estimator and the pooled model as bridges are added (Section~\ref{sec:res-within}).

\begin{figure}[h]
\centering
\widefig{figures/fig3_tradeoff.pdf}
\caption{The within-site estimator keeps the reverse-set AUC above chance from the smallest number of bridges, whereas the pooled model approaches it only as $L$ falls. Within-site test AUC against reverse-set AUC as bridges are added (ABIDE-I, CC200; means over $50$ repeats). Each point is one number of bridges per site $b$; $b=0$ is the class-pure baseline, for which L3 is undefined. Red: pooled model C0; blue: within-site estimator L3.}
\label{fig:tradeoff}
\end{figure}

\section{Deep within-site objectives}\label{app:deep}

We asked whether the within-site principle carries over to a deep representation, and how it compares with deep domain generalization. All deep models used the BrainNetCNN backbone~\cite{kawahara2017} of the robustness check.

\emph{Frozen encoder with linear heads.} A denoising autoencoder ($20\%$ symmetric masking, Adam with learning rate $10^{-4}$, weight decay $5\times10^{-4}$, batch size $14$, $100$ epochs) was trained on an unlabeled pool made of the training participants, without their labels, and all ABIDE-I participants at sites that were neither training nor test sites of the repeat; input standardization was fitted on the same pool. Test-site participants and the reverse set never entered it. Three linear heads were fitted on the frozen embedding of the training participants: logistic regression (E-CE), linear Group DRO with one group per site (E-GDRO), and L3 (E-L3).

\emph{End-to-end baselines.} From random initialization, with the settings of the robustness check (SGD, learning rate $10^{-4}$, momentum $0.9$, weight decay $5\times10^{-4}$, batch size $14$, $50$ epochs, no tuning): ERM (D-ERM), online Group DRO with one group per site (D-GDRO), and a multi-source domain-adversarial network~\cite{ganin2016} (D-DANN) with a site discriminator behind a gradient-reversal layer, which also saw the same unlabeled pool. In timing runs, the D-DANN loss diverged because the encoder could enlarge its unnormalized features without bound; before the formal runs, and judged on training losses only, we normalized the discriminator input with a layer normalization without affine parameters.

The four cells were $k\in\{2,4\}$, $b=5$, one-way and two-way, with $50$ repeats; in the $k=2$ two-way cell, repeats $0$--$9$ had been used in a pilot and only repeats $10$--$49$ were used for the criteria. Table~\ref{tab:deep} shows the results. With the same embedding, the L3 head kept the reverse-set AUC at $0.58$--$0.60$, against $0.34$--$0.51$ for the logistic head, at similar test AUC. Criterion D1 was not met because the difference in the $k=4$ two-way cell was $0.07$, below the threshold of $0.10$. The end-to-end baselines had reverse-set AUCs of $0.22$--$0.51$. In this D-DANN configuration, the site remained linearly decodable from the representation. After the final epoch, a logistic regression on the raw $128$-dimensional layer, with five folds stratified by site within the training and external pool, identified the site with an accuracy of $0.69$--$0.91$ (chance $0.08$; all $50$ repeats per cell), against $0.42$--$0.48$ for ERM. This probe measures decodability within a pool the encoder had seen, not transfer to unseen sites. The discriminator of D-DANN, whose accuracy during the last epoch was near $0$, read the layer-normalized representation, so the two accuracies are not directly comparable. All deep models stayed below the linear models in test AUC ($0.55$--$0.58$ against $0.58$--$0.63$).

\begin{table}[h]
\centering
\caption{Deep models (ABIDE-I, CC200, $b=5$). Within-site test AUC / reverse-set AUC, means over repeats ($40$ for $k=2$ two-way, $50$ otherwise). Linear C0 and L3 are the models of the main text.}
\label{tab:deep}
\small
\resizebox{\linewidth}{!}{%
\begin{tabular}{lcccc}
\toprule
 & $k=2$, one-way & $k=2$, two-way & $k=4$, one-way & $k=4$, two-way\\
\midrule
E-CE (frozen encoder) & $0.567$ / $0.341$ & $0.579$ / $0.406$ & $0.570$ / $0.471$ & $0.573$ / $0.513$\\
E-GDRO (frozen encoder) & $0.564$ / $0.358$ & $0.576$ / $0.400$ & $0.569$ / $0.497$ & $0.573$ / $0.515$\\
E-L3 (frozen encoder) & $0.554$ / $0.577$ & $0.584$ / $0.598$ & $0.569$ / $0.592$ & $0.583$ / $0.586$\\
D-ERM & $0.554$ / $0.268$ & $0.574$ / $0.354$ & $0.568$ / $0.407$ & $0.577$ / $0.513$\\
D-GDRO & $0.554$ / $0.267$ & $0.573$ / $0.355$ & $0.567$ / $0.405$ & $0.576$ / $0.513$\\
D-DANN & $0.554$ / $0.218$ & $0.565$ / $0.300$ & $0.563$ / $0.379$ & $0.571$ / $0.492$\\
\midrule
Linear C0 & $0.596$ / $0.199$ & $0.616$ / $0.315$ & $0.607$ / $0.388$ & $0.623$ / $0.524$\\
Linear L3 & $0.582$ / $0.603$ & $0.619$ / $0.642$ & $0.604$ / $0.622$ & $0.632$ / $0.647$\\
\bottomrule
\end{tabular}}
\end{table}

\section{Robustness, statistical tests and further details}\label{app:robust}

\subsection{Robustness variants}

Table~\ref{tab:robust} gives the three pre-specified robustness criteria for the three ABIDE-I variants. The BrainNetCNN cost at $k=1$ only just exceeded its threshold of $0.03$, and the bridging comparison rested on $5$ to $16$ repeats in which both designs were feasible.

\begin{table}[h]
\centering
\caption{Robustness variants of ABIDE-I. Q1: cost of purity M $-$ P at $k=1$ and $k=6$, and reverse-set AUC of P at $k=1$ and $k=6$. Q2: $E_{\mathrm{subj}}-$ truth at $k=2$, and mean absolute error of $E_{\mathrm{subj}}$ and LSPO at $k=2$. Q3: C0 $-$ S0 in test and reverse-set AUC, two-way $b=5$, $k=2$, with the number of feasible repeats.}
\label{tab:robust}
\small
\resizebox{\linewidth}{!}{%
\begin{tabular}{lccc}
\toprule
 & QC subset ($n=873$) & Harvard-Oxford & BrainNetCNN\\
\midrule
Q1 cost, $k=1\to6$ & $0.073\to0.018$ & $0.064\to0.016$ & $0.032\to0.020$\\
Q1 reverse-set AUC, $k=1\to6$ & $0.02\to0.34$ & $0.03\to0.37$ & $0.11\to0.37$\\
Q2 bias of $E_{\mathrm{subj}}$ & $+0.37$ & $+0.38$ & $+0.35$\\
Q2 MAE, $E_{\mathrm{subj}}$ / LSPO & $0.37$ / $0.14$ & $0.38$ / $0.15$ & $0.35$ / $0.14$\\
Q3 test / reverse-set difference & $+0.012$ / $+0.18$ (15) & $+0.031$ / $+0.22$ (16) & $+0.035$ / $+0.16$ (5)\\
\bottomrule
\end{tabular}}
\end{table}

\subsection{ComBat null test}

Before ComBat was applied to the bridged designs, we ran it on pure noise with the dimensions of each atlas: standard normal features plus one $N(0,1)$ shift vector per site, no diagnostic effect, $k=2$, $b=5$, one-way and two-way, five repeats each. The pre-specified requirement was a mean reverse-set AUC within $[0.40,0.60]$. With empirical Bayes, ComBat gave $0.000$ in all four groups. Empirical Bayes shrinks each site's location estimate towards a common prior, so a fraction of each site shift remains after harmonization; in an unbalanced design this remainder is correlated with the class, and across thousands of features a linear classifier exploits it. In a small simulation, exact removal of site means and ComBat without empirical Bayes both gave about $0.5$. Without empirical Bayes, three of the four groups passed ($0.53$, $0.41$, $0.44$; Harvard-Oxford two-way $0.34$). We therefore report only the test AUC of ComBat, which was within $0.01$ of C0 for both versions in all four $b=5$ cells and both atlases.

\subsection{Corrected resampled tests}

Table~\ref{tab:tests} lists all $82$ corrected resampled $t$-tests~\cite{nadeau2003}, grouped into the six families within which Holm correction was applied. For each paired contrast it gives the mean difference over the $J$ repeats, the number of repeats with a positive difference ($n_{>0}$), and the two-sided $p$-values before and after Holm correction. One test remains below $0.05$ after correction. As explained in Section~\ref{sec:metrics}, the correction is an approximation in this design; the counts of positive repeats describe stability across overlapping splits and are not independent trials. Cells: $k$, $b$, one-way (1w) or two-way (2w).

{\footnotesize\setlength{\tabcolsep}{2.5pt}
\begin{longtable}{lllcrrrrr}
\caption{Corrected resampled $t$-tests of the paired contrasts (test: within-site AUC on new sites; reverse: reverse-set AUC).}\label{tab:tests}\\
\toprule
Cohort & Contrast & Metric & Cell & $J$ & Mean & $n_{>0}$ & $p$ & $p_{\mathrm{Holm}}$\\
\midrule
\endfirsthead
\toprule
Cohort & Contrast & Metric & Cell & $J$ & Mean & $n_{>0}$ & $p$ & $p_{\mathrm{Holm}}$\\
\midrule
\endhead
\bottomrule
\endfoot
\multicolumn{9}{l}{\textit{1. Cost of class purity}}\\
ABIDE-I & M $-$ P & test & k1 & 50 & $+0.057$ & 37 & 0.674 & 1.000\\
ABIDE-I & M $-$ P & test & k2 & 50 & $+0.053$ & 40 & 0.615 & 1.000\\
ABIDE-I & M $-$ P & test & k4 & 50 & $+0.028$ & 34 & 0.734 & 1.000\\
ABIDE-I & M $-$ P & test & k6 & 50 & $+0.015$ & 29 & 0.868 & 1.000\\
\multicolumn{9}{l}{\textit{2. Evaluation}}\\
ABIDE-I & $E_{\mathrm{subj}}-$ truth & test & k2 & 50 & $+0.365$ & 50 & 0.006 & 0.035\\
ABIDE-I & abs.\ error LSPO $-$ subj. & test & k2 & 50 & $-0.204$ & 6 & 0.371 & 1.000\\
ABIDE-I & $E_{\mathrm{subj}}-$ truth & test & k4 & 50 & $+0.269$ & 50 & 0.062 & 0.309\\
ABIDE-I & abs.\ error LSPO $-$ subj. & test & k4 & 50 & $-0.154$ & 4 & 0.419 & 1.000\\
ABIDE-I & $E_{\mathrm{subj}}-$ truth & test & k6 & 50 & $+0.189$ & 50 & 0.107 & 0.428\\
ABIDE-I & abs.\ error LSPO $-$ subj. & test & k6 & 50 & $-0.100$ & 7 & 0.536 & 1.000\\
\multicolumn{9}{l}{\textit{3. Bridging}}\\
ABIDE-I & C0 $-$ B0 & reverse & k2/b5/1w & 50 & $+0.079$ & 50 & 0.086 & 1.000\\
ABIDE-I & C0 $-$ B0 & test & k2/b5/1w & 50 & $+0.022$ & 39 & 0.634 & 1.000\\
ABIDE-I & C0 $-$ B0 & reverse & k2/b5/2w & 50 & $+0.183$ & 50 & 0.046 & 0.642\\
ABIDE-I & C0 $-$ B0 & test & k2/b5/2w & 50 & $+0.040$ & 44 & 0.471 & 1.000\\
ABIDE-I & C0 $-$ B0 & reverse & k4/b5/1w & 50 & $+0.117$ & 50 & 0.159 & 1.000\\
ABIDE-I & C0 $-$ B0 & test & k4/b5/1w & 50 & $+0.025$ & 39 & 0.543 & 1.000\\
ABIDE-I & C0 $-$ B0 & reverse & k4/b5/2w & 50 & $+0.253$ & 50 & 0.017 & 0.273\\
ABIDE-I & C0 $-$ B0 & test & k4/b5/2w & 50 & $+0.042$ & 43 & 0.415 & 1.000\\
ABIDE-I & C0 $-$ S0 & reverse & k2/b5/1w & 29 & $+0.096$ & 29 & 0.085 & 1.000\\
ABIDE-I & C0 $-$ S0 & test & k2/b5/1w & 29 & $+0.015$ & 18 & 0.740 & 1.000\\
ABIDE-I & C0 $-$ S0 & reverse & k2/b5/2w & 16 & $+0.238$ & 16 & 0.021 & 0.312\\
ABIDE-I & C0 $-$ S0 & test & k2/b5/2w & 16 & $+0.044$ & 15 & 0.451 & 1.000\\
ABIDE-I & C0 $-$ S0 & reverse & k4/b5/1w & 35 & $+0.133$ & 35 & 0.107 & 1.000\\
ABIDE-I & C0 $-$ S0 & test & k4/b5/1w & 35 & $+0.024$ & 28 & 0.504 & 1.000\\
ABIDE-I & C0 $-$ S0 & reverse & k4/b5/2w & 15 & $+0.291$ & 15 & 0.048 & 0.642\\
ABIDE-I & C0 $-$ S0 & test & k4/b5/2w & 15 & $+0.028$ & 11 & 0.564 & 1.000\\
\multicolumn{9}{l}{\textit{4. Within-site estimation}}\\
ABIDE-I & L3 $-$ C0 & reverse & k2/b2/1w & 50 & $+0.437$ & 50 & 0.026 & 0.546\\
ABIDE-I & L3 $-$ C0 & test & k2/b2/1w & 50 & $-0.029$ & 15 & 0.769 & 1.000\\
ABIDE-I & L3 $-$ C0 & reverse & k2/b2/2w & 50 & $+0.425$ & 50 & 0.026 & 0.546\\
ABIDE-I & L3 $-$ C0 & test & k2/b2/2w & 50 & $-0.003$ & 27 & 0.973 & 1.000\\
ABIDE-I & L3 $-$ C0 & reverse & k2/b5/1w & 50 & $+0.404$ & 50 & 0.017 & 0.386\\
ABIDE-I & L3 $-$ C0 & test & k2/b5/1w & 50 & $-0.014$ & 18 & 0.846 & 1.000\\
ABIDE-I & L3 $-$ C0 & reverse & k2/b5/2w & 50 & $+0.337$ & 50 & 0.014 & 0.341\\
ABIDE-I & L3 $-$ C0 & test & k2/b5/2w & 50 & $+0.005$ & 28 & 0.919 & 1.000\\
ABIDE-I & L3 $-$ C0 & reverse & k2/b10/1w & 50 & $+0.373$ & 50 & 0.024 & 0.517\\
ABIDE-I & L3 $-$ C0 & test & k2/b10/1w & 50 & $-0.004$ & 24 & 0.931 & 1.000\\
ABIDE-I & L3 $-$ C0 & reverse & k2/b10/2w & 50 & $+0.184$ & 50 & 0.042 & 0.801\\
ABIDE-I & L3 $-$ C0 & test & k2/b10/2w & 50 & $+0.007$ & 33 & 0.746 & 1.000\\
ABIDE-I & L3 $-$ C0 & reverse & k4/b2/1w & 50 & $+0.277$ & 50 & 0.135 & 1.000\\
ABIDE-I & L3 $-$ C0 & test & k4/b2/1w & 50 & $-0.021$ & 19 & 0.801 & 1.000\\
ABIDE-I & L3 $-$ C0 & reverse & k4/b2/2w & 50 & $+0.238$ & 50 & 0.098 & 1.000\\
ABIDE-I & L3 $-$ C0 & test & k4/b2/2w & 50 & $-0.001$ & 26 & 0.990 & 1.000\\
ABIDE-I & L3 $-$ C0 & reverse & k4/b5/1w & 50 & $+0.234$ & 50 & 0.138 & 1.000\\
ABIDE-I & L3 $-$ C0 & test & k4/b5/1w & 50 & $-0.003$ & 26 & 0.961 & 1.000\\
ABIDE-I & L3 $-$ C0 & reverse & k4/b5/2w & 50 & $+0.123$ & 50 & 0.160 & 1.000\\
ABIDE-I & L3 $-$ C0 & test & k4/b5/2w & 50 & $+0.009$ & 35 & 0.725 & 1.000\\
ABIDE-I & L3 $-$ C0 & reverse & k4/b10/1w & 50 & $+0.186$ & 50 & 0.117 & 1.000\\
ABIDE-I & L3 $-$ C0 & test & k4/b10/1w & 50 & $+0.002$ & 27 & 0.958 & 1.000\\
ABIDE-I & L3 $-$ C0 & reverse & k4/b10/2w & 50 & $-0.011$ & 20 & 0.811 & 1.000\\
ABIDE-I & L3 $-$ C0 & test & k4/b10/2w & 50 & $+0.007$ & 38 & 0.551 & 1.000\\
\multicolumn{9}{l}{\textit{5. Methods on class-pure data}}\\
ABIDE-I & LSPO tuning $-$ B0 & reverse & k2 & 50 & $+0.004$ & 26 & 0.786 & 1.000\\
ABIDE-I & LSPO tuning $-$ B0 & test & k2 & 50 & $+0.001$ & 19 & 0.925 & 1.000\\
ABIDE-I & LSPO tuning $-$ B0 & reverse & k4 & 50 & $+0.003$ & 21 & 0.868 & 1.000\\
ABIDE-I & LSPO tuning $-$ B0 & test & k4 & 50 & $-0.000$ & 16 & 0.982 & 1.000\\
ABIDE-I & LSPO tuning $-$ B0 & reverse & k6 & 50 & $+0.003$ & 20 & 0.864 & 1.000\\
ABIDE-I & LSPO tuning $-$ B0 & test & k6 & 50 & $+0.000$ & 18 & 0.954 & 1.000\\
ABIDE-I & NAP $-$ B0 & reverse & k2 & 50 & $+0.002$ & 29 & 0.866 & 1.000\\
ABIDE-I & NAP $-$ B0 & test & k2 & 50 & $-0.001$ & 24 & 0.959 & 1.000\\
ABIDE-I & NAP $-$ B0 & reverse & k4 & 50 & $-0.001$ & 21 & 0.961 & 1.000\\
ABIDE-I & NAP $-$ B0 & test & k4 & 50 & $-0.000$ & 25 & 0.983 & 1.000\\
ABIDE-I & NAP $-$ B0 & reverse & k6 & 50 & $+0.001$ & 27 & 0.958 & 1.000\\
ABIDE-I & NAP $-$ B0 & test & k6 & 50 & $-0.001$ & 24 & 0.947 & 1.000\\
ABIDE-I & Group DRO $-$ ERM & reverse & k2 & 50 & $-0.003$ & 12 & 0.757 & 1.000\\
ABIDE-I & Group DRO $-$ ERM & test & k2 & 50 & $+0.000$ & 25 & 0.963 & 1.000\\
ABIDE-I & Group DRO $-$ ERM & reverse & k4 & 50 & $-0.001$ & 18 & 0.943 & 1.000\\
ABIDE-I & Group DRO $-$ ERM & test & k4 & 50 & $+0.000$ & 27 & 0.958 & 1.000\\
\multicolumn{9}{l}{\textit{6. Replication}}\\
ABIDE-II & $E_{\mathrm{subj}}-$ truth & test & k2 & 50 & $+0.376$ & 50 & 0.024 & 0.383\\
ABIDE-II & abs.\ error LSPO $-$ subj. & test & k2 & 50 & $-0.184$ & 5 & 0.474 & 1.000\\
ABIDE-II & abs.\ error LSPO $-$ subj. & test & k4 & 50 & $-0.102$ & 9 & 0.727 & 1.000\\
ABIDE-II & M $-$ P & test & k1 & 50 & $+0.050$ & 35 & 0.728 & 1.000\\
ABIDE-II & cost, $k$4 $-$ $k$1 & test & k4-k1 & 50 & $-0.050$ & 19 & 0.819 & 1.000\\
ABIDE-II & P, $k$4 $-$ $k$1 & reverse & k4-k1 & 50 & $+0.237$ & 50 & 0.192 & 1.000\\
ABIDE-II & L3 $-$ C0 & reverse & k2/b5/1w & 50 & $+0.391$ & 50 & 0.036 & 0.533\\
ABIDE-II & L3 $-$ C0 & test & k2/b5/1w & 50 & $-0.023$ & 17 & 0.799 & 1.000\\
ABIDE-II & L3 $-$ C0 & reverse & k2/b5/2w & 50 & $+0.239$ & 50 & 0.055 & 0.766\\
ABIDE-II & L3 $-$ C0 & test & k2/b5/2w & 50 & $-0.006$ & 21 & 0.901 & 1.000\\
ADHD-200 & C0 $-$ S0 & reverse & k1/b5/2w & 48 & $+0.076$ & 46 & 0.462 & 1.000\\
ADHD-200 & C0 $-$ S0 & test & k1/b5/2w & 48 & $+0.013$ & 31 & 0.900 & 1.000\\
ADHD-200 & L3 $-$ C0 & reverse & k1/b5/1w & 60 & $+0.518$ & 60 & 0.099 & 1.000\\
ADHD-200 & L3 $-$ C0 & test & k1/b5/1w & 60 & $+0.005$ & 34 & 0.967 & 1.000\\
ADHD-200 & L3 $-$ C0 & reverse & k1/b5/2w & 60 & $+0.440$ & 60 & 0.125 & 1.000\\
ADHD-200 & L3 $-$ C0 & test & k1/b5/2w & 60 & $+0.004$ & 35 & 0.954 & 1.000\\
\end{longtable}

}

\subsection{Published multi-site methods}

Table~\ref{tab:published} gives all results of the two published methods. In the original transductive setting the reverse set entered training without labels; in the second setting it was kept out and scored after training. The mixed design M has no reverse set, so it was run once. Subject-level cross-validation uses class-pure training data only and is the same for both settings. The graph network with the CC200 atlas was run in the transductive setting only.

\begin{table}[h]
\centering
\caption{Published multi-site methods on ABIDE-I ($n=40$, $50$ repeats). Within-site test AUC and reverse-set AUC for $k=2$ / $k=4$; $E_{\mathrm{subj}}$: subject-level five-fold cross-validation of B0 at $k=2$. Trans.: original transductive setting; out: reverse set kept out of training.}
\label{tab:published}
\small\setlength{\tabcolsep}{3.5pt}
\resizebox{\linewidth}{!}{%
\begin{tabular}{llcccccc}
\toprule
Method & Setting & B0 test & B0 reverse & C0 test & C0 reverse & M test & $E_{\mathrm{subj}}$\\
\midrule
Logistic regression & inductive & $.575$ / $.581$ & $.120$ / $.269$ & $.614$ / $.623$ & $.301$ / $.524$ & $.627$ / $.609$ & $.940$\\
Graph network, HO & trans. & $.589$ / $.634$ & $.107$ / $.543$ & $.620$ / $.642$ & $.608$ / $.619$ & $.604$ / $.605$ & $1.000$\\
Graph network, HO & out & $.559$ / $.591$ & $.017$ / $.222$ & $.599$ / $.630$ & $.306$ / $.557$ & & \\
Graph network, CC200 & trans. & $.598$ / $.627$ & $.070$ / $.524$ & $.643$ / $.641$ & $.646$ / $.638$ & $.618$ / $.605$ & $1.000$\\
MIDA pipeline & trans. & $.625$ / $.642$ & $.544$ / $.576$ & $.643$ / $.648$ & $.596$ / $.622$ & $.612$ / $.611$ & $.352$\\
MIDA pipeline & out & $.519$ / $.527$ & $.347$ / $.404$ & $.598$ / $.625$ & $.429$ / $.582$ & & \\
\bottomrule
\end{tabular}}
\end{table}


\bibliographystyle{elsarticle-num}
\bibliography{refs}
\end{document}
